% ============================================================
% CHAPTER 4: PROJECT OVERVIEW
% MCDA 5585 / 5586 Major Project Report
% Bhavik Kantilal Bhagat | A00494758 | Saint Mary's University
% ============================================================

\chapter{Project Overview}
\label{chap:project}

This chapter provides the production context that established the
monitoring gap and motivated the work, then describes the four project
streams and their engineering mandates.

% ────────────────────────────────────────────────────────────
\section{Project Background and the Motivating Incident}
\label{sec:project_background}
% ────────────────────────────────────────────────────────────

\subsection{Pre-Internship Monitoring Posture}

Before my internship began in mid-Mar~2026, the IA-SRE team operated
without a centralized observability stack. Monitoring relied entirely on
an \emph{ad~hoc} model: business teams in Dublin, London, and Halifax
would notice anomalies in fund processing outputs and raise tickets. There
were no CloudWatch alarms, no dashboards aggregating health across the six
platforms, and no mechanism for the SRE team to receive proactive
notification of degrading conditions --- the team was blind to
\emph{saturation} and \emph{error rate} signals until they had already
caused downstream business impact.

\subsection{CAIS-Pricing Silent Outage (Jan--Feb 2026)}
\label{sec:cais_silent_outage}

Between \textbf{Jan~20 and Feb~17, 2026}, the CAIS-Pricing Extract GenAI
pipeline under the \texttt{rpacfs1} platform stopped delivering output for
\textbf{29 consecutive days} without a single alert firing. The root cause
was gradual exhaustion of the external CTI AI model API token, which
controls the async SQS callback that triggers the downstream Lambda
stages; without that callback, three downstream functions dropped to zero
invocations while two upstream functions kept running, making the pipeline
appear healthy. The outage was detected only when business users reported
missing pricing data --- 29~days after onset. The post-mortem (IISS-430)
identified the absence of throughput alarms, SQS queue depth monitoring,
and end-to-end output verification as the root causes of the failure going
unnoticed. Sub-task IISS-435 was assigned directly to me as the immediate
remediation and became the starting point for the full monitoring programme.

\subsection{The Meridian Incident (Apr~2026)}

The Meridian Incident is described in full detail in
Section~\ref{sec:meridian_platform} of Chapter~\ref{chap:about_company}.
In brief: a surge of automated system events from Citco Recs overwhelmed
Meridian's OpenSearch consumption layer, causing Kafka consumer lag to
reach approximately \textbf{9~million messages} and halting daily
statement processing for 7.5~hours. The SRE team received no alert; the
Dublin business team escalated over three and a half hours after the
first failure signal. The post-mortem recorded: \textit{``No proactive
monitoring/alerting in place on the SRE side for Meridian event volume
anomalies.''} This became the primary driver for IISS-487.

\subsection{From Incident to Programme}

Both incidents shared the same root cause: no proactive monitoring. The
post-mortems established two specific gaps --- no Kafka consumer lag alarm,
and no centralized visibility across the IA-IT-SRE Infrastructure --- that
directly motivated the three work streams described in the following
section. Kishor elevated IISS-487 to \textbf{Critical} priority on
May~6, 2026, and my internship pivoted from exploratory dashboard research
to a full-scale observability platform build.

% ────────────────────────────────────────────────────────────
\section{Project Overview}
\label{sec:project_overview}
% ────────────────────────────────────────────────────────────

The two motivating incidents established a clear engineering mandate: build
proactive observability across the entire IA-IT-SRE Infrastructure. Three
work streams were pursued sequentially and in parallel, each one directly
motivated by the findings and limitations of the previous. A fourth
ongoing responsibility --- RPA production support --- ran throughout the
internship period.

\subsection{Work Stream 1: CloudWatch Dashboards and IA Infrastructure Monitoring}

Following the CAIS-Pricing Silent Outage, the immediate
remediation was to build a unified CloudWatch monitoring dashboard starting
with Lambda --- the AWS service at the centre of that incident. The goal was
to surface invocation count, error rate, duration, and throttle metrics
for every Lambda function across all six platforms, giving the SRE team
visibility that had not existed before.

Once Lambda was covered, the scope was extended incrementally to the other
AWS services in the IA-IT-SRE Infrastructure: SQS queues (queue depth and
oldest message age), EC2 instances (CPU and memory via CloudWatch Agent),
ECS clusters and tasks (via CloudWatch Container Insights), EKS node metrics,
and database services. Each resource type was added as a separate CloudWatch
dashboard, establishing the tag-based resource discovery patterns ---
including the \texttt{AppName} fallback strategy for shared \texttt{BudgetCode}
environments --- that all later work built upon.

These CloudWatch dashboards were effective for per-service visibility, but
they exposed a fundamental limitation: CloudWatch dashboards are siloed
per AWS service and do not support dynamic cross-service filtering. There was
no way to view Lambda, SQS, and ECS health for a single platform side by side
in one workspace. This limitation directly motivated Work Stream~2.
Key deliverables: IISS-435, IISS-455, IISS-461, IISS-469, IISS-482--486.

\subsection{Work Stream 2: Grafana AMG Centralized Dashboard}

Following the Meridian Incident, Kishor elevated IISS-487 to
\textbf{Critical} priority. The goal was to build a proof-of-concept (PoC)
centralized observability platform using \textbf{Amazon Managed Grafana
(AMG)} --- addressing what the per-service CloudWatch dashboards could not:
a unified, cross-platform, cross-service view of the entire IA-IT-SRE
Infrastructure in a single workspace.

The scope was defined by the Meridian post-mortem action items and refined
through requirements sessions with Kishor: the dashboard needed to cover
every AWS service already monitored in Work Stream~1 (Lambda, SQS, MSK/Kafka,
EC2, ECS, EKS, ECR, RDS Aurora, DocumentDB, Redis/Valkey, Neptune DB, OpenSearch),
support dynamic filtering by platform and AWS service via template variables,
and include CloudWatch alarms with SNS notifications to the SRE team. The
broader research objective was to evaluate whether this AWS-native stack could
achieve Dynatrace-equivalent observability --- documenting the feature parity
achieved, the gaps that remain, the limitations encountered, and the
incremental cost at each tier.

The dashboard was built on CloudWatch as the sole data source (via the native
Grafana CloudWatch plugin), using Metric Insights SQL for cross-resource
aggregation and CloudWatch Logs Insights for log-derived panels. An
Infrastructure-as-Code provisioner --- a Lambda-backed CloudFormation pipeline
--- was developed to dynamically discover the AWS service fleet at deploy time
and synthesize per-service alarm definitions, enabling a single deployment to
cover all monitored AWS services without manual enumeration. Key limitations
discovered and documented during this work stream --- including the Metric
Insights SQL percentile restriction, the SEARCH() visibility gap for idle
resources, and the 64~KB Secrets Manager payload limit --- are detailed in
Chapter~\ref{chap:methodologies}. Key deliverable: IISS-487 (140 hours).

\subsection{Work Stream 3: IA Infrastructure Inventory API}

A critical limitation surfaced during the Grafana dashboard build: the
\texttt{WHERE} clause in CloudWatch Metric Insights SQL does not support
pattern-based filtering (no wildcard or regex on dimension values). To filter
a panel to a specific platform, an exact dimension value match was required ---
but there was no reliable way to enumerate the available values dynamically from
within Grafana. The \texttt{dimension\_values()} function returned partial
results for large fleets, and cross-platform queries required hardcoding platform
names into template variable queries.

To solve this, the Lambda discovery script that had been built for the
CloudWatch dashboards was extended into a full-fledged \textbf{Flask REST API}:
the \texttt{ia-it-sre-infrastructure-inventory} service. The API discovers AWS
services across all six platforms, regions, and environments in real time at
query time, and exposes structured resource lists via REST endpoints. A Grafana
dashboard panel can call the API to retrieve the current list of Lambda
functions, ECS clusters, or SQS queues for a given platform, and use the
response to populate template variable options dynamically --- ensuring the
dashboard always reflects the live fleet without requiring manual updates.

The API uses a discoverer registry pattern: a registry maps
\texttt{(resource\_type, platform)} tuples to concrete discoverer classes,
allowing new platforms or AWS service types to be onboarded by implementing
a single class. The service was built with Python and \texttt{flask[async]},
using \texttt{boto3} for all AWS API calls and Pydantic for response validation.
The test suite achieved 99\% code coverage across unit, integration, and
property-based tests. Key deliverable: IISS-507.

\subsection{Work Stream 4: RPA Production Support}

In parallel with the three engineering work streams, I participated in
dedicated RPA production support rotation weeks as part of the team's
\textbf{Follow the Sun} support model. The North America slot covers
11~AM -- 7~PM~EST, and I held this slot once every four weeks, serving
as the primary point of contact for all RPA and automation-related
incidents during the shift.

The support model is structured across three tiers. \textbf{Level~1} is
handled by the AI/RPA Support Agent (24/7). \textbf{Level~2} is owned by
the Innovation IT Support team (us) email channel, operating under
a Follow the Sun model with a 1-hour acknowledgement SLA and 8-hour
resolution target. \textbf{Level~3} is the SRE team
(\texttt{ia\_it\_sre@citco.com}), engaged for complex production issues,
root cause analysis, and service improvement, with a 1-day response SLA.
As an SRE intern I operated across both L2 and L3 responsibilities during
support rotations.

At the start of each shift I reviewed all open items across the relevant
channels: the \textbf{SD (Service Desk) Portal} for incoming tickets, the
\textbf{innovation\_it\_support@citco.com} email inbox for L2 escalations,
the \textbf{IA\_INCIDENTS@citco.com} distribution list for incident
notifications, and \textbf{Microsoft Teams} messages and channel
notifications from business teams and the outgoing shift. I also reviewed
the Confluence shift handover log left by the Manila (Phillipinse) shift to 
pick up any carried-over items or approaching SLA breaches. Ticket ownership rules
required that whoever acknowledged a ticket owned it to resolution ---
handover notes had to include actions taken, pending items, blockers, and
clear next steps for the incoming shift.

Incident work followed the full lifecycle. I monitored running Blue Prism
and UiPath bot queues (particularly the VPL queue opening at 5~PM~EST
daily), triaged service desk tickets, and for each issue performed root
cause analysis: investigating CloudWatch logs, UiPath Orchestrator
heartbeat status, Blue Prism control room logs, and application-level
error messages. Where a fix was within support scope --- updating a
Corporate Actions FileMask for client onboarding, retriggering a stuck
MESO pipeline, correcting a configuration file --- I resolved it directly.
Where the root cause required a code change beyond support scope, I
prepared a documented RCA, proposed a resolution approach, and raised or
assigned a JIRA to the appropriate development team, copying \texttt{ia\_it\_sre@citco.com} per
the incident escalation procedure.

Notable incidents handled during support rotations included the CAIS
Deadline credential failure (IISS-706), where I identified the Secrets
Manager rotation gap and delivered a fail-fast hotfix; the MESO month-end
stuck tasks (IISS-741, Critical) requiring manual pipeline retrigger; and
the VPL Pricing GTL incorrect booking (IISS-876), where I performed the
RCA and coordinated the fix with the business team.
